\begin{papercard}{malki2024}
\cf{Stages} \stagetag{A} \stagetag{D} \stagetag{W} \stagetag{C} \stagetag{I}. Closest end-to-end comparator for the thesis: store screening, structured store disclosures, privacy policies and a privacy-focused interface inspection, triangulated across surfaces.

\cf{Aim and research questions} Exploratory study of the privacy practices of popular female mHealth apps as communicated through Data safety sections, privacy policies and app interfaces \pc{malki2024}{2}{as communicated through Data safety sections}. RQ1 asks for declared practices on collection, storage, sharing, deletion and safeguarding; RQ2 for available in-app privacy features (deletion, portability, consent/opt-out) and their usability; RQ3 for inconsistencies between declarations and actual in-app mechanisms and their safety implications after Roe v.\ Wade \pc{malki2024}{2}{Do inconsistencies exist between the privacy practices de}.

\cf{Design} Mixed methods: thematic analysis of Data safety sections and policies plus a hybrid usability inspection combining cognitive walkthrough and heuristic evaluation, then triangulation of the three surfaces to find disconnects \pc{malki2024}{4}{We then triangulated our findings to iden}. Ethics: a minimal-risk registration form; the King's College London ethics office judged that no formal review was needed \pc{malki2024}{4}{After filling out a minimal}.

\cf{Sampling frame and search} Google Play only (Android $>$70\,\% of devices); apps had to be available in \emph{both} the US and UK stores \pc{malki2024}{4}{We included apps which were available in both the US and}. Search terms were derived from prior work and refined iteratively by testing each term in the Play Store search for variety of relevant results \pc{malki2024}{4}{improved our terms by testing each with the Play Store search}. Final 12 terms: \emph{female health, women's health, reproductive health, fertility, ovulation, menstrual, menstruation, period, contraception, birth control, pregnancy, pregnant} \pc{malki2024}{4}{We used the following finalised set of terms for our search}. A Python script collected all results in March 2023, giving 184 unique apps \pc{malki2024}{4}{This returned a total of 184 unique apps}.

\cf{Eligibility and selection} Sequential exclusions with counts: zero user ratings ($n=54$), not English ($n=3$), not updated for over a year ($n=35$), no direct female reproductive-health service ($n=29$) \pc{malki2024}{4}{We excluded apps which had zero user ratings}. The remaining 63 were screened by store screenshots and the top 10--20 user reviews; apps without at least one tracking or predictive feature ($n=7$) and non-consumer apps ($n=4$) were removed \pc{malki2024}{4}{scanning the top 10-20 user reviews}. This left 52 eligible apps (40 period/fertility trackers, 12 pregnancy/baby apps) \pc{malki2024}{4}{This left a total of 52 eligible apps}. Ranking signal: \emph{number of reviews in the preceding 12 months}, chosen so that newer but rising apps were not disadvantaged; the top 20 were analysed \pc{malki2024}{4}{We ranked these apps based on the number of reviews}. Table A.1 records, per app, genre, developer, developer type, release date, downloads, ratings, 2023 reviews and version \pc{malki2024}{20}{Details of female mHealth apps selected for analysis}.

\cf{Materials captured} (1) Data safety section fields: data types collected and shared, purposes, encryption in transit, deletion option; these were extracted manually, and the same fields were searched in the policy by keyword for comparison \pc{malki2024}{4}{manually extracted this information from the Data safety section}. (2) Privacy policies: 19 unique policies for 20 apps (one developer shared a generic policy), 18 linked in Google Play and two found on the developer website \pc{malki2024}{5}{We identified a total of 19 unique privacy policies across the 20}. (3) App interfaces, inspected on an Android device.

\cf{Coding scheme} Deductive start from prior FemTech privacy work and GDPR data-subject rights, chosen as a comprehensive and widely used yardstick \pc{malki2024}{4}{We began by drafting an initial codebook based on prior}. Eleven top-level categories: policy scope; policy version; user data types; collection methods; purposes; storage/retention; transfers; safeguarding measures; regional legislation; user privacy rights; language and readability \pc{malki2024}{4}{Our codebook included the following key categories}. The full codebook (Table C.1) gives a definition and a real example quotation per code \pc{malki2024}{22}{Full codebook used in our analysis of female mHealth app privacy policies}. Examples: \emph{Limitations} of safeguards, i.e.\ transferring responsibility to the user \pc{malki2024}{23}{Mentions of any limitations of the developers}; \emph{Consequences of exercising rights} \pc{malki2024}{23}{Any unavoidable outcomes or consequences of the user exer}; \emph{Ambiguous}, \emph{Contradictory} (also against Data safety sections and app features) and \emph{Complex language} \pc{malki2024}{23}{Information that is conflicting or contradictory}.

\cf{Codebook development and coders} Iterative team procedure \pc{malki2024}{4}{codebook underwent multiple rounds of iteration}: the first author piloted the deductive draft on five random policies; each of the five authors coded a \emph{different} policy and noted uncaptured text; the team revised (more purposes and third-party sharing); then three authors repeatedly coded sets of 3--5 policies independently and discussed disagreements until code saturation \pc{malki2024}{4}{three authors conducted iterations of independently}. Final coding: the first author re-coded all policies line-by-line in batches of five in MAXQDA, with team meetings after each batch \pc{malki2024}{4}{all app policies were re-coded line-by-line by the first}. \textbf{No inter-rater reliability was calculated} \pc{malki2024}{4}{We did not calculate inter-rater reliability (IRR)}.

\cf{Walkthrough procedure} A \emph{heuristic walkthrough} (hybrid of cognitive walkthrough and heuristic evaluation) with a preparation stage (target user and tasks) and an analysis stage (simulated use guided by Nielsen's heuristics and privacy-by-design principles) \pc{malki2024}{5}{We conducted a heuristic walkthrough}. The target user was a first-time or novice user tracking their own sexual and reproductive health on Android \pc{malki2024}{5}{as a first-time or novice user of a female mHealth app}. All 20 apps were installed and every screen, button and feature was inspected. Use cases were condensed into common tasks, and tasks that access, input or manage personal data became the focus \pc{malki2024}{5}{Tasks which involved accessing}. Analysis: the first and last authors each inspected every app independently. Apps were wiped and re-installed to simulate first use \pc{malki2024}{5}{After wiping and re-installing the 20 apps}. Each task was walked step by step, asking whether the user would know how to proceed, identify the correct action and interpret the feedback \pc{malki2024}{5}{whether the user would know how to proceed with each task}. Each task was also checked against Nielsen's ten heuristics and eight privacy-by-design guidelines (Appendix B) \pc{malki2024}{21}{Summary of Privacy-by-Design guidelines}. Problems were logged with possible fixes, collated in Miro and aggregated by constant comparison \pc{malki2024}{5}{Common design problems were identified through constant}.

\cf{Key findings (method-relevant)} Data safety sections and policies agreed on data types and purposes but diverged on sharing, encryption and deletion. All eight apps that declared no sharing described third-party sharing in the policy \pc{malki2024}{5}{All 8/20 apps which claimed in their Data safety sections}. Policy length averaged 4{,}453 words (SD 2{,}636); the shortest covered only website cookies \pc{malki2024}{6}{The shortest app policy was just 300 words}. Scope problems: one policy covered several apps and services \pc{malki2024}{6}{Policy Scope. We analysed language relating to the scope}. Versioning: 12/19 had been updated within 12 months and four had no revision date \pc{malki2024}{6}{12/19 privacy policies had been updated within the last}. The UI inspection produced a typology of three deletion mechanisms: account only, separate account/data deletion, and full deletion \pc{malki2024}{10}{we identified three overarching techniques for}. Only 7/20 apps had granular privacy controls \pc{malki2024}{10}{Only 7/20 apps implemented granular pri}. Interface practices were also observed: a greyed-out skip option and a disguised optional surname field \pc{malki2024}{10}{which is a typical indication that an option is disabled}. Exports were incomplete or behind a paywall \pc{malki2024}{12}{Only 8/20 apps provided an in-app mecha}. Coined \emph{cascading consent} for consent that propagates to ever more third-party contexts \pc{malki2024}{14}{a phenomenon to which we refer to as cascading}.

\cf{Limitations stated} Google Play only, so iOS differences are untested \pc{malki2024}{16}{we only included apps available in Google Play in our analysis}. Manual policy coding is not scaled. The codebook could support automation and readability metrics \pc{malki2024}{16}{could upscale our approach}. The inspection was done by the research team without users, so it does not mirror long-term everyday use \pc{malki2024}{16}{our usability inspection was conducted by the re}. No legal specialists were involved \pc{malki2024}{16}{Lastly, legal specialists were not formally involved in the research}.

\cf{Use for this thesis} Directly reusable: the counted exclusion cascade, recent-review ranking as a popularity signal with an explicit rationale, the cross-surface comparison fields (Data safety vs.\ policy vs.\ UI), the code definition + example format, the target-user/task preparation and the three CW questions. Not transferable as-is: the US/UK English market, the Android-only frame, and the coding design. Their iterative team coding was followed by single-coder final coding without IRR, so it is \emph{not} the frozen-manual, two-coder reliability design required in the thesis.
\end{papercard}
